% Clase 24 — la ecuación de Laplace discretizada: en cada punto interior de la
% grilla de paso a, el potencial es el promedio de sus cuatro vecinos. A la
% derecha, el problema del ejemplo: caja con phi = 0 en el borde y una placa
% interior a phi_0.
\begin{tikzpicture}[scale=1]

  % ---------- estencil ----------
  \foreach \xx in {0,1,...,4} { \draw[figgray!50, line width=0.4pt] (\xx,0) -- (\xx,4); }
  \foreach \yy in {0,1,...,4} { \draw[figgray!50, line width=0.4pt] (0,\yy) -- (4,\yy); }
  \foreach \xx in {0,1,...,4} { \foreach \yy in {0,1,...,4} { \fill[figgray] (\xx,\yy) circle (1.3pt); } }
  \fill[figred] (2,2) circle (3pt);
  \foreach \px/\py in {3/2, 1/2, 2/3, 2/1} {
    \fill[figblue] (\px,\py) circle (2.6pt);
    \draw[figvecaux] (\px,\py) -- ($(\px,\py)!0.62!(2,2)$);
  }
  \node[figlbl, text=figred, fill=white, inner sep=1pt] at (2.45,1.62) {$\phi(x,y)$};
  \node[figlbl, anchor=west] at (3.08,2.25) {$\phi(x+a,y)$};
  \draw[figgray, line width=0.6pt, {Latex[length=1.5mm]}-{Latex[length=1.5mm]}] (0,-0.3) -- (1,-0.3);
  \node[figlbl, anchor=north] at (0.5,-0.35) {$a$};
  \node[figlbl, anchor=north] at (2,-0.9) {(a) promedio de los vecinos};

  % ---------- problema ----------
  \begin{scope}[xshift=6.2cm]
    \fill[figgray!10] (0,0) rectangle (4,4);
    \draw[figblue, line width=1.3pt] (0,0) rectangle (4,4);
    \fill[figred!30] (1.3,2.2) rectangle (2.6,2.8);
    \draw[figred, line width=1pt] (1.3,2.2) rectangle (2.6,2.8);
    \node[figlbl] at (1.95,2.5) {$\phi_0$};
    \node[figlbl, fill=figgray!10, inner sep=1pt] at (2.0,0.35) {$\phi = 0$ en el borde};
    \node[figlbl] at (2.0,1.3) {$\nabla^2\phi = 0$};
    \node[figlbl, anchor=north] at (2,-0.9) {(b) el problema};
  \end{scope}

\end{tikzpicture}
